\documentclass{article}%
\usepackage[T1]{fontenc}%
\usepackage[utf8]{inputenc}%
\usepackage{lmodern}%
\usepackage{textcomp}%
\usepackage{lastpage}%
\usepackage[head=40pt,margin=0.5in,bottom=0.6in]{geometry}%
\usepackage{graphicx}%
%
\title{\textbf{Docentes protestaron en Portuguesa por bajos salarios}}%
\author{El Nacional Web}%
\date{22/11/2018}%
%
\begin{document}%
\normalsize%
\maketitle%
\textbf{URL: }%
http://www.el{-}nacional.com/noticias/protestas/docentes{-}protestaron{-}portuguesa{-}por{-}bajos{-}salarios\_260711\newline%
%
\textbf{Periodico: }%
EN, %
ID: %
260711, %
Seccion: %
Protestas\newline%
%
\textbf{Palabras Claves: }%
Educación, Nicolás Maduro, Economía, Protestas, Gobierno, Portuguesa\newline%
%
\textbf{Derecho: }%
2.2%
, Otros Derechos: %
2.3%
, Sub Derechos: %
2.2.1, 2.3.4%
\newline%
%
\textbf{EP: }%
SI\newline%
\newline%
%
\textbf{\textit{El gremio de educadores indicó que el dinero no les alcanza para cubrir sus gastos}}%
\newline%
\newline%
%
\includegraphics[width=300px]{263.jpg}%
\newline%
%
El gremio de docentes del~municipio Turén, ~estado Portuguesa, protestó~este jueves~en rechazo a los bajos sueldos que reciben por su labor.%
\newline%
%
Los manifestantes indicaron que el dinero no les alcanza para cubrir sus gastos. Aseguran que solo les alcanza para comprar un cartón de huevos.%
\newline%
%
Los educadores gritaron consignas y sostuvieron~pitos y pancartas. “Sueldo de docentes para el presidente”, es parte de lo que vociferan.%
\newline%
%
Desde hace algunos meses, se han registrado protestas de educadores en varias entidades del país.~ Exigen~un salario justo al presidente Nicolás Maduro.%
\newline%
%
\end{document}